%% Guidage en rotation par deux roulements : arrêts axiaux de chaque bague.
%% La règle de montage est la réponse : elle passe par \rappel.
\documentclass[tikz,border=4pt]{standalone}
\usepackage{liaisons}
\begin{document}
\begin{tikzpicture}[x=1cm,y=1cm,
  axe/.style={line width=0.5pt, draw=black!55, dash pattern=on 9pt off 3pt on 2pt off 3pt},
  repere/.style={line width=0.5pt, draw=black!60},
  etiq/.style={font=\small, text=black!80, inner sep=2pt},
  arret/.style={-{Stealth[length=5pt,width=4.5pt]}, line width=1pt, draw=solideE}]

%% ---------------- l'arbre, épaulé à gauche ----------------
\hachures[draw=solideA, line width=0.9pt]{
  (0,0) -- (0,0.62) -- (1.5,0.62) -- (1.5,0.95) -- (2.4,0.95) -- (2.4,0.62)
  -- (8.2,0.62) -- (8.2,0) -- cycle}
\node[font=\small\bfseries, text=solideA, anchor=east] at (-0.15,0.31) {Arbre};

%% ---------------- les deux roulements ----------------
\foreach \x in {2.9, 6.6} {
  \hachures[draw=solideA, line width=0.8pt]{(\x,0.62) rectangle (\x+1.0,1.02)}
  \draw[line width=0.8pt, draw=black!70, fill=black!12] (\x+0.5,1.28) circle (0.24);
  \hachuresb[draw=solideB, line width=0.8pt]{(\x,1.54) rectangle (\x+1.0,1.94)}
}

%% ---------------- le logement ----------------
\hachuresb[draw=solideB, line width=0.9pt]{
  (1.1,1.94) -- (1.1,2.6) -- (8.2,2.6) -- (8.2,1.94) -- (7.6,1.94) -- (7.6,1.54)
  -- (7.0,1.54) -- (7.0,1.94) -- (3.9,1.94) -- (3.9,1.54) -- (3.3,1.54) -- (3.3,1.94) -- cycle}
\node[font=\small\bfseries, text=solideB, anchor=west] at (8.35,2.27) {Logement};

\draw[axe] (-0.5,0) -- (8.8,0);

%% ---------------- les arrêts axiaux ----------------
\draw[arret] (2.55,0.78) -- (2.85,0.78);
\node[etiq, text=solideE, anchor=south east, font=\footnotesize] at (2.9,0.95) {épaulement};
\draw[line width=2pt, draw=solideE, line cap=round] (7.72,0.66) -- (7.72,1.0);
\draw[arret] (7.6,0.82) -- (7.9,0.82);
\node[etiq, text=solideE, anchor=south west, font=\footnotesize] at (7.7,1.02) {anneau élastique};

%% ---------------- la règle : c'est la réponse ----------------
\node[anchor=north west, font=\small, align=left, inner sep=3pt] at (-0.5,-0.4) {%
  \rappel{chaque bague est arrêtée \textbf{des deux côtés} : sinon l'arbre se promène dans son logement ;\\
  la bague qui tourne par rapport à la charge est montée \textbf{serrée}, l'autre \textbf{glissante}}};
\end{tikzpicture}
\end{document}
